\section*{2. Literature review and conceptual framework}
\addcontentsline{toc}{section}{2. Literature review and conceptual framework}

\subsection*{2.1 From biological plausibility to demonstrated association}
\addcontentsline{toc}{subsection}{2.1 From biological plausibility to demonstrated association}

The hypothesised pathway from WASH to child growth is biologically direct: inadequate water and sanitation increase exposure to faecal--oral pathogens, raising the frequency of diarrhoeal disease and sub-clinical environmental enteric dysfunction, which together redirect energy and nutrients away from linear growth \parencite{CummingCairncross2016}. This mechanism is not challenged even by the strongest recent evidence against basic WASH interventions reducing stunting: the authors of the three major WASH trials discussed below state explicitly that "the biological plausibility of WASH as public health interventions is not challenged by these findings" \parencite[consensus message 2]{Cumming2019}, while arguing that the \textit{specific, basic} interventions they tested were insufficient to translate that plausible mechanism into a detectable growth effect. Black et al. (\citeyear{Black2013}) place WASH-related infection among the proximate determinants of child undernutrition in their widely cited \textit{Lancet} framework, alongside inadequate diet and maternal and child care practices, and Victora et al. (\citeyear{Victora2010}) document that growth faltering is concentrated in the window from birth to roughly 24 months. What this literature review is concerned with is a narrower, empirical question: does a \textit{facility-type classification} --- whether a household has an "improved" water source or sanitation facility, in the standard WHO/UNICEF JMP sense used throughout this thesis --- show a detectable, policy-relevant statistical association with child growth once the obvious confounders are accounted for, as distinct from the question of whether WASH-related infection can plausibly affect growth at all.

\subsection*{2.2 Observational evidence and the confounding problem}
\addcontentsline{toc}{subsection}{2.2 Observational evidence and the confounding problem}

A substantial observational literature finds WASH variables associated with child anthropometric outcomes in DHS and related survey data. Rakotomanana et al. (\citeyear{Rakotomanana2020}) link WHO/UNICEF JMP water and sanitation indicators to linear growth in children aged 6--23 months across several East African DHS surveys and find associations in the expected direction. Addae et al. (\citeyear{Addae2024}), using Ghana's most recent national survey, find household water, sanitation, and socioeconomic status jointly associated with stunting and wasting among children aged 6--23 months. Günther and Fink (\citeyear{Gunther2010}) pool 172 DHS surveys and estimate that improved water and sanitation are associated with lower child mortality, and Spears (\citeyear{Spears2013}) documents that cross-country variation in average child height correlates with cross-country variation in open-defecation rates.

The central methodological difficulty common to this literature, and squarely faced in this thesis, is that household and community WASH status are not randomly assigned: both are strongly correlated with household wealth, parental education, and local infrastructure, all of which independently predict child nutrition through channels unrelated to pathogen exposure. Headey and Palloni (\citeyear{HeadeyPalloni2019}) address this directly, constructing a panel of 442 subnational regions across 59 countries from repeated Demographic and Health Surveys and estimating difference-in-differences regressions that exploit within-region change over time, net of region fixed effects, rather than cross-sectional variation alone. Their results are mixed across outcomes and should be read precisely rather than summarised as a single attenuation story: improved water access is statistically insignificantly associated with most outcomes, although water piped directly into the home specifically predicts reductions in child stunting; sanitation improvements predict sizeable reductions in diarrhoea prevalence and child mortality but are \textit{not} associated with changes in stunting or wasting once the within-region panel design is used. This last finding --- a null sanitation--stunting association once region-level time-invariant confounding is netted out --- is a closer and more precise parallel to this thesis's own null household- and community-sanitation findings (Section 5) than a general "attenuation" story would suggest, and it is drawn on that way in Section 6 rather than as evidence that WASH associations with child health "generally" weaken under better controls, which overstates what this one study, on its own, establishes.

\subsection*{2.3 Credible causal evidence: what the major trials actually found, and what they did not test}
\addcontentsline{toc}{subsection}{2.3 Credible causal evidence: what the major trials actually found, and what they did not test}

Randomised evidence on WASH and child growth comes primarily from three cluster-randomised controlled trials, described precisely here, trial by trial and arm by arm, because loose paraphrase of their design risks both mischaracterising a genuinely different exposure structure from the cross-sectional literature in Section 2.2 and 2.5, and attributing a finding from one arm to a different intervention than the one that produced it. WASH-Benefits Bangladesh, WASH-Benefits Kenya, and the Sanitation Hygiene Infant Nutrition Efficacy (SHINE) trial in Zimbabwe each enrolled pregnant women and their children \textit{in utero} and followed them prospectively for 18--24 months, assigning interventions near the start of that follow-up window, not at a single arbitrary survey date unrelated to each child's own age. This is a materially different exposure-timing structure from the DHS-based observational studies discussed in Section 2.2 and 2.5, and the exposure-timing limitation described in Section 1.2 for this thesis's own baseline design does not apply to these three trials.

Each trial used a factorial design with \textbf{separate arms}, not a single combined treatment: a water arm (point-of-use chlorination), a sanitation arm (an improved household or compound-level latrine), a hygiene arm (a handwashing station with soap), a combined-WASH arm (all three together), a nutrition arm (lipid-based nutrient supplementation and feeding-behaviour promotion, with no WASH component), and a combined WASH-plus-nutrition arm, each compared against a control. On linear growth specifically, \textbf{every individual WASH-related arm --- water, sanitation, hygiene, and combined WASH --- showed no statistically detectable effect in all three trials} \parencite[pp. 2, 5 and Fig. 2]{Cumming2019}. The nutrition arm, which contained no WASH component, showed a separate, modest positive effect on linear growth, described by the trial authors as a small gain despite high compliance \parencite[p. 6]{Cumming2019}; the combined WASH-plus-nutrition arm showed \textbf{no additional growth benefit over the nutrition arm alone} \parencite[p. 2]{Cumming2019} --- that is, adding WASH to nutrition did not improve on nutrition by itself. This is reported here, trial-component by trial-component, specifically so that the nutrition arm's modest growth effect is not read as a WASH finding: it was produced by an arm with no WASH intervention at all. On diarrhoea, results were genuinely mixed across the WASH-related arms and sites --- a substantial relative reduction in Bangladesh (against a low baseline prevalence), no detectable effect in Kenya, and no detectable effect in Zimbabwe \parencite[pp. 2--5, Fig. 2]{Cumming2019}.

"No statistically detectable effect" is the precise description used throughout this discussion, not "proven to be exactly zero": these are adequately powered trials with high internal validity, and the authors themselves describe the growth result as consistent evidence across three sites, but a null result in a hypothesis test establishes an absence of detected effect at the trial's power and sample size, not mathematical certainty that the true effect is zero. The authors' own interpretation, discussed next, is also explicit that this finding should not be generalised to \textit{every possible} WASH intervention.

The consensus statement summarising these trials is explicit that this finding does not challenge the biological plausibility of WASH affecting growth (Section 2.1) and attributes the null growth result instead to the limited scope of the specific interventions tested: household-level, "basic" WASH services (in the JMP ladder sense) that left community-level open defecation and the broader water-supply infrastructure largely unchanged, rather than the "safely managed," piped, networked services the authors argue may be needed for a detectable population-level effect \parencite{Cumming2019}. The authors explicitly call for research on more comprehensive, "transformative" WASH packages rather than concluding that WASH cannot matter for growth. This distinction matters directly for how this thesis's own results should be read: the exposure measured throughout this thesis --- household- or community-level "improved" facility status --- is, like the trial interventions, a facility-type classification rather than a comprehensive intervention package, so a null or weak association in this thesis's data is consistent with, though does not on its own confirm, the trial literature's finding that facility-level WASH alone is often insufficient to move linear growth; it is equally consistent with confounding masking a real effect, with measurement error, or with this thesis's own specific limitations (Section 6). This thesis does not interpret its null or weak findings as evidence that WASH investment in general is ineffective.

A related, quasi-experimental literature targets infant mortality rather than anthropometric growth, using a different identification strategy worth describing precisely. Geruso and Spears (\citeyear{GerusoSpears2015}) investigate why Muslim infants in India survive to their first birthday at a substantially higher rate than Hindu infants, despite Muslim households having lower wealth, consumption, and education on average --- a puzzle documented by prior literature and not explained by household characteristics alone. The authors show that neighbourhoods with a higher share of Muslim households are, on most observable dimensions, \textit{worse} places to raise a child, consistent with Muslims' general socioeconomic disadvantage in India --- with one specific, striking exception: open defecation. Hindus practise open defecation at a rate roughly 25 percentage points higher than Muslims nationally, a religious behavioural difference not explained by relative wealth, so neighbourhoods with more Muslim households have systematically better sanitation conditions even though they are worse on nearly every other observable predictor of child health. This lets the authors isolate a sanitation-specific channel from general neighbourhood quality, because the two move in \textit{opposite} directions across religious composition rather than together, and they find this sanitation channel accounts for a substantial share of the Muslim mortality advantage. This thesis cites the working-paper version of this study specifically (version and verification details recorded in \texttt{bibliography\_verification\_ledger.md}, not restated here); it is cited for its identification strategy --- exploiting a source of neighbourhood sanitation variation plausibly distinct from general local development --- as an example of the kind of design this thesis's own four-country, cross-sectional comparison cannot replicate, not as direct evidence about child growth, which Geruso and Spears do not study.

\subsection*{2.4 Household versus community exposure}
\addcontentsline{toc}{subsection}{2.4 Household versus community exposure}

Most of the DHS-based observational studies reviewed here measure WASH at the household level only. A smaller, more recent strand examines community-level or spatially aggregated WASH exposure as a distinct construct. Donohue, Church, Assaf, and Mayala (\citeyear{Donohue2023}), in a DHS Spatial Analysis Report, construct community-level improved-sanitation coverage (from a geospatial model, not a leave-one-out DHS-census rate) and relate it to childhood stunting in two countries, Nigeria and Zambia (both 2018 DHS), using multilevel logistic regression. The result is mixed, not a general finding, and is stated precisely here rather than paraphrased into one: in Nigeria, community coverage remained significantly associated with lower odds of stunting after full adjustment, though only modestly (adjusted OR = 0.97, 95\% CI [0.95, 0.99] per 10-percentage-point increase), while household-level improved sanitation was no longer significant in the same adjusted model; in Zambia, community coverage was not significantly associated with stunting in either the unadjusted or adjusted model, and neither was household sanitation. This motivates constructing and reporting both household and community exposures as a worthwhile empirical question --- the two measures can diverge --- but it is not evidence that community coverage generally "carries explanatory relevance beyond" household facilities, which overstates a result significant in one of the two countries studied. Skoufias and Vinha (\citeyear{SkoufiasVinha2026}) develop a related community-externality framework for child health in Sub-Saharan Africa using DHS data from twenty countries, and their result is similarly specific rather than general: they find \textbf{no significant externality from community-level water coverage}, but do find community-level improved sanitation significantly associated with lower fever/diarrhoea incidence and higher height-for-age, with effects varying between urban and rural areas. This is a closer and more precise parallel to this thesis's own baseline finding (Section 5.2) that community water and sanitation do not behave identically under full adjustment, though neither source's result should be read as confirming this thesis's specific coefficients, which come from a different sample, country set, and model. This thesis's leave-one-out community construction follows this same logic; the author is not aware of a study pairing this exact leave-one-out construction with a household-level comparison in the same four-country sample, though this thesis does not claim to have conducted an exhaustive search of the literature sufficient to assert this combination is unprecedented more broadly --- only that reviewing household and community WASH side by side, under the fixed-effects structure each exposure's own construction requires, is the specific design choice this thesis makes and did not find replicated in the sources reviewed here. Balk et al. (\citeyear{Balk2004}), working on child mortality in West Africa with spatially referenced survey data, is an earlier methodological precedent for treating DHS-cluster-level geography as a unit of analysis distinct from the household, though applied to mortality rather than anthropometric outcomes.

A specific identification issue arises when community WASH is measured this way: because the leave-one-out rate is itself a property of the cluster rather than of the individual child, attempting to estimate it with cluster fixed effects --- the natural approach for household WASH --- would remove nearly all of its substantive variation, leaving a coefficient identified almost entirely by the small mechanical residual that arises because a different household is excluded from each child's own calculation (Section 4 derives this precisely). This thesis therefore uses region, not cluster, fixed effects for the community specification.

\subsection*{2.5 Exposure timing in single-round, cross-sectional survey designs specifically}
\addcontentsline{toc}{subsection}{2.5 Exposure timing in single-round, cross-sectional survey designs specifically}

The limitation introduced in Section 1.2 applies to studies that measure WASH status once, at a single survey date, and relate it to the anthropometric status of children of varying current ages within that same round --- a description that fits this thesis's own baseline design and the DHS cross-sectional studies reviewed in Section 2.2 \parencite{Rakotomanana2020,Addae2024,Gunther2010,Spears2013}, but does \textbf{not} describe the major WASH trials (Section 2.3), which assign and measure exposure prospectively from near each child's own birth, nor Headey and Palloni's (\citeyear{HeadeyPalloni2019}) panel design, which exploits change over successive survey rounds within the same subnational regions rather than a single round's cross-section. Age controls and age-interaction terms --- which this thesis also estimates for its baseline sample (Section 4, Section 5) --- describe whether the \textit{current} survey-date association differs by how old a child is today; they do not reconstruct what a child's own environment was during their own early life, and this thesis is explicit that these are different empirical questions. This single-round limitation, specific to this design class, is the reason this thesis's historical extension attempts a different, early-life-window exposure assignment within its own separate sample (Section 1.2), not a reason to expect the two samples' results to isolate an exposure-timing effect when compared to each other (Section 6.2).

\subsection*{2.6 Geospatial exposure assignment, and a matching-method precedent from outside the WASH literature}
\addcontentsline{toc}{subsection}{2.6 Geospatial exposure assignment, and a matching-method precedent from outside the WASH literature}

Assigning a continuous, gridded exposure surface to a survey respondent's location requires handling two well-documented DHS features. First, coordinates are randomly displaced for confidentiality \parencite{Burgert2013,PerezHeydrich2013}, which affects any extraction method: whichever pixel or pixels are read, they are read around the cluster's \textit{recorded}, displaced coordinates, not its true location, and no extraction method recovers the true location or corrects for the (undisclosed, per-cluster) displacement actually applied. Second, a value still has to be read from the surface at that location, and more than one method exists for doing so --- a containing-pixel extraction, which reads the single pixel the point falls in, or a buffer-based extraction, which averages pixels within some radius, by simple area weighting or by population weighting if population data for the area are available. This thesis uses an area-weighted buffer average as its primary method and a containing-pixel extraction as a labelled sensitivity (Section 3.6; Section 5.8); the two are not interchangeable descriptions of the same procedure, and the choice between them is made explicit here rather than treated as if only one were possible. Johnson, Jacob, and Brown (\citeyear{JohnsonJacobBrown2013}) provide an earlier precedent for linking satellite-derived environmental surfaces (forest cover) to DHS cluster locations to study a different child-health outcome. Blom, Ortiz-Bobea, and Hoddinott (\citeyear{Blom2022}) match historical, geocoded weather data to each child's own developmental window across repeated DHS rounds --- a design this thesis's historical extension borrows structurally. This source was reviewed at the level of its published abstract and design description, not independently verified claim-by-claim in this pass (recorded in \texttt{bibliography\_verification\_ledger.md}); it is cited here, and only here, for its general matching-method structure --- matching a time-varying, geocoded exposure to each child's own developmental window rather than to the survey date --- not for any specific numerical finding, and not as evidence about WASH or nutrition, which that study does not examine.

\subsection*{2.7 This thesis's contribution and the gap it addresses}
\addcontentsline{toc}{subsection}{2.7 This thesis's contribution and the gap it addresses}

This thesis combines three design elements that this literature review has not found combined in one study: household and leave-one-out community WASH compared directly within the same sample and controls (Section 2.4); and, in a separate, bounded, and explicitly preliminary extension, an attempt to assign modelled local WASH coverage to each child's own early-life window using geospatial matching (Section 2.6), within the historical sample's own three countries. That a specific four-country combination has not been studied before is not, on its own, evidence of a substantive methodological contribution, and this thesis does not rely on country selection as its claim to contribution. The methodological contribution claimed is narrower: a direct, consistently controlled comparison of the two exposure levels, and a transparent, fully documented attempt --- including its own unresolved issues --- at an early-life exposure-matching design that neither Donohue et al. (\citeyear{Donohue2023}) nor Skoufias and Vinha (\citeyear{SkoufiasVinha2026}) implement; both examine current-round community WASH exposure (geospatially modelled or DHS-census-based respectively), not a time-varying exposure matched to each child's own developmental window, and neither source is cited here as calling for that specific extension --- the gap is this thesis's own observation, not an attributed recommendation. The gap this leaves open, and does not claim to close, is the same one Headey and Palloni (\citeyear{HeadeyPalloni2019}) and the trial literature \parencite{Cumming2019} already identify at a larger scale: credibly separating a WASH-specific effect from the broader package of local development improved facilities tend to accompany. This thesis's historical extension narrows the exposure-timing dimension of that problem within its own sample; it does not resolve the confounding dimension, and it does not resolve the exposure-timing dimension for the baseline sample, since the two samples differ in several other respects simultaneously (Section 6.2).
